\section{Methods}

\subsection{Inference-only baselines}

The first family contains no weight updates. \texttt{baseline1} directly prompts the local base model with the Tangut string. \texttt{baseline2} augments the prompt with dictionary glosses. \texttt{baseline2\_1\_cot} further restructures the prompt into a three-step process intended to encourage alignment, literal drafting, and rewriting. This family is closest in spirit to lexicon-augmented or terminology-aware MT, except that the lexicon is injected through prompting rather than through decoding constraints or model-side features \cite{chen2017incorporating, hokamp2017lexically}.

We also add a deliberately stronger prompt-only comparison based on DeepSeek-V3.2 with few-shot dictionary prompting. This baseline is served through OpenRouter, receives the same dictionary gloss construction, uses a much stricter title-only system prompt, and is given four fixed development-set few-shot exemplars. The prompt explicitly instructs the model to prefer compact title form, avoid explanatory paraphrase, resist blindly copying generic sutra names from dictionary entries, and emit only the final title string. Decoding uses temperature 0.0 with a 512-token cap. The goal is not to build a straw baseline but to answer a reviewer-style question directly: if a modern frontier Chinese model with carefully engineered prompting already solves the task, then claims about the necessity of local adaptation should be narrowed accordingly.

These methods are important because they show what can be achieved from symbol injection alone. In the current results, local dictionary prompting is substantially better than zero-shot generation, but it often expands compact titles into explanatory paraphrases rather than reconstructing the reference form. The frontier DeepSeek variant shows that prompt-only inference becomes much stronger once both the model and the prompt design are upgraded aggressively.

\subsection{Synthetic SFT baselines}

The second family uses synthetic training data. Synthetic examples are built from monolingual classical Chinese strings and a Tangut dictionary by first constructing a Chinese-character-to-Tangut-character map from dictionary explanations and then replacing characters according to one of four regimes.
\begin{itemize}
\item \texttt{baseline3} keeps a mixed-script input by replacing only 30--70\% of replaceable Chinese characters with Tangut.
\item \texttt{baseline3\_1\_unk} enforces pure Tangut inputs by replacing every dictionary miss with \texttt{[UNK]}.
\item \texttt{baseline3\_3\_semantic} also enforces pure inputs, but replaces misses through vector-space semantic projection instead of \texttt{[UNK]}.
\item \texttt{baseline3\_2\_multitask} keeps the same pure-input construction as the \texttt{UNK} branch but changes the target format to include explicit alignment scaffolding.
\end{itemize}

The multitask target is concrete rather than notional:
\begin{quote}
\small\ttfamily
<dict\_match>\{Tangut input\}</dict\_match>\\
<literal>\{character-level gloss proxy\}</literal>\\
\{final Chinese title\}
\end{quote}
This gives the model both a copy-style alignment channel and a final title-generation target in the same sequence. The central design question is therefore whether the model benefits more from input purity, approximate semantic preservation, or explicit alignment structure. At a high level, this family plays the role that synthetic parallel data often plays in low-resource MT: it converts monolingual or weakly aligned resources into supervised signals that a translation model can absorb \cite{sennrich2016monolingual}. The current evidence favors explicit alignment structure.

\subsection{Preference optimization}

The final family applies DPO to an SFT base using reward-defined chosen/rejected pairs \cite{rafailov2023dpo}. Candidate preferences are generated by sampling five outputs per training prompt from the SFT model at temperature 0.8 and top-$p$ 0.95, scoring them with
\[
r(y) = \mathrm{lex}(y) - 0.01 \log \mathrm{ppl}(y),
\]
where $\mathrm{lex}(y)$ is the repository's lexical-coverage score. We then keep the top-scoring output as chosen and the bottom-scoring output as rejected whenever the reward gap is at least 0.05. We study this family in two stages because the repository now contains both a failure case and a repaired variant.

The legacy pipeline first selects the best SFT base among the mixed, \texttt{UNK}, and semantic variants using chrF++, then constructs preference pairs from that base. In the current snapshot, that selector chooses the \texttt{UNK} model, yielding the legacy checkpoint \texttt{final\_v2}. This run is still important because it exposes the failure mode directly: if the base model is weaker and the reward is partly misaligned, DPO can degrade title reconstruction rather than improve it.

The follow-up DPO runs start instead from the strongest supervised base, the multitask SFT model. After auditing the preference data, we remove invalid rows and keep only pairs whose reward gap exceeds a threshold, testing both a looser filter ($\ge 0.2$, 1,996 pairs) and a stricter filter ($\ge 0.4$, 498 pairs). For each filtered set we train two variants: standard sigmoid DPO and a noise-aware alternative that combines TRL's robust DPO loss with WPO-style reweighting for off-policy preferences \cite{chen2024robustdpo,pang2024wpo}. This turns the DPO family into a controlled diagnosis of three factors at once: base-model choice, pair-quality filtering, and loss robustness.

\subsection{Training and judge details}

All learned models start from Qwen2.5-7B-Instruct and use LoRA adapters. The SFT runs use three epochs, learning rate $2\times10^{-4}$, batch size 4, gradient accumulation 4, rank-64 LoRA, and maximum sequence length 512. The DPO runs use two epochs, learning rate $5\times10^{-5}$, batch size 1, gradient accumulation 8, $\beta=0.1$, rank-32 LoRA, cosine decay, and maximum sequence length 512. The robust variants add label smoothing 0.1 and WPO-style weighting. The full run configurations are saved alongside each checkpoint.

The reference-aware judge is a separate supplementary evaluator, not a training signal. It uses an Azure-hosted GPT-5.4 deployment and scores each candidate on four integer dimensions: reference agreement, source faithfulness, title-style fitness, and overall quality. The judge prompt sees the Tangut source, the gold reference, and the candidate simultaneously and explicitly penalizes expansions of short titles into explanatory modern sentences. We use this judge because the repository's original reference-free judge is not strong enough to anchor the task, but we still treat it as supplemental evidence rather than a replacement for human evaluation.
